\documentclass[11pt]{article}
\usepackage[a4paper,margin=2.2cm]{geometry}
\usepackage{xcolor}
\usepackage{fontspec}
\usepackage{xeCJK}
\usepackage{enumitem}
\setmainfont{Times New Roman}
\setCJKmainfont{SimSun}
\setlist[itemize]{leftmargin=1.4em,itemsep=0.35em,topsep=0.35em}
\setlength{\parindent}{0pt}
\setlength{\parskip}{0.55em}
\pagestyle{plain}
\begin{document}
{\color[HTML]{1F5FD2}\LARGE\bfseries 自闭症儿童教育支持策略\par}
{\color[HTML]{667085}\small 星轨原创教育资源：用于家庭与学校共同制定支持计划。\par}
\vspace{0.8em}
{\color[HTML]{111827}\large\bfseries 理解支持需求\par}
\begin{itemize}
\item 自闭症儿童可能在沟通、社交理解、感觉调节和活动转换方面需要明确、稳定且可预期的支持。
\item 选择策略前先观察行为发生前后的线索。同样的行为可能来自焦虑、感觉过载、语言不清或缺少表达工具。
\end{itemize}
{\color[HTML]{111827}\large\bfseries 课堂策略\par}
\begin{itemize}
\item 使用视觉日程，提前预告转换，并为每个活动提供清晰的开始和结束信号。
\item 把任务拆成小步骤，儿童完成一步就给予即时、具体的反馈，而不是等整项任务结束再评价。
\item 根据需要提供安静角、降噪耳机、活动休息或手部小工具，帮助儿童恢复参与状态。
\end{itemize}
{\color[HTML]{111827}\large\bfseries 沟通支持\par}
\begin{itemize}
\item 接纳口语、手势、图片、书写和辅助沟通系统等多种表达方式，不把说话作为唯一合格的沟通形式。
\item 成人应先示范如何表达请求、拒绝、评论和求助，再逐步等待儿童尝试。
\end{itemize}
\vfill
{\color[HTML]{667085}\footnotesize 本材料为应用内原创教育内容，不能替代医疗、法律或正式评估建议。\par}
\end{document}
